%% UrbanAI'26 submission draft.
%% Full research paper track (8-10 pages incl. references).
%%
%% Every number in this document is produced by the scripts in sidewalk/ and
%% was verified against the data on 2026-08-05. Regenerate before submission:
%% the city deployments are live and counts drift.
%%
%% For review, the CFP does not state whether the process is anonymous. If it
%% is, switch to \documentclass[sigconf,anonymous]{acmart} and drop the author
%% block. Confirm with the organisers before submitting.

\documentclass[sigconf]{acmart}

\AtBeginDocument{%
  \providecommand\BibTeX{{Bib\TeX}}}

\setcopyright{acmlicensed}
\copyrightyear{2026}
\acmYear{2026}
\acmDOI{XXXXXXX.XXXXXXX}
\acmConference[UrbanAI '26]{The 4th ACM SIGSPATIAL International Workshop on
  Advances in Urban AI}{November 3, 2026}{Riverside, CA, USA}
\acmISBN{978-1-4503-XXXX-X/26/11}

\begin{document}

\title{What Quality Control Removes: Contested Labels in Crowdsourced
Urban Accessibility Data}

\author{Himson Chapagain}
\affiliation{%
  \institution{Texas State University}
  \city{San Marcos}
  \state{Texas}
  \country{USA}}
\email{rnx16@txstate.edu}

\begin{abstract}
Cities increasingly rely on crowdsourced assessments of the built environment to
decide where to spend accessibility budgets. Those assessments are cleaned
before use, and cleaning treats reviewer disagreement as noise. Across
811{,}368 labels from twelve Project Sidewalk deployments, we show that the
platform's published quality criterion retains 54.5\% of reviewed labels overall
but only 8.6\% of labels reviewers split on and 6.8\% of labels they were
unsure about, discarding over 90\% of both in all twelve cities. The removal is
spatially uneven, varying three- to eight-fold beyond sampling noise across
neighbourhoods and surviving controls for what gets labelled where. It is also
consequential for downstream models: training a validity classifier on the
filtered record makes it assert more than 90\% confidence on 48\% of contested
labels, against 30\% when the ambiguous cases are retained at matched sample
size, with no gain in discriminative accuracy. The effect replicates across
three model families. We argue that contested labels should be routed rather
than deleted, and that the two kinds of disagreement require different
destinations: better imagery for one, a policy decision for the other.
\end{abstract}

\begin{CCSXML}
<ccs2012>
<concept><concept_id>10003120.10003130.10003233</concept_id>
<concept_desc>Human-centered computing~Collaborative and social computing systems and tools</concept_desc>
<concept_significance>500</concept_significance></concept>
<concept><concept_id>10002951.10003227.10003351</concept_id>
<concept_desc>Information systems~Data mining</concept_desc>
<concept_significance>300</concept_significance></concept>
</ccs2012>
\end{CCSXML}
\ccsdesc[500]{Human-centered computing~Collaborative and social computing systems and tools}
\ccsdesc[300]{Information systems~Data mining}

\keywords{crowdsourcing, data quality, urban accessibility, annotator
disagreement, model calibration, volunteered geographic information}

\maketitle

\section{Introduction}

Whether a stretch of pavement is passable in a wheelchair is not always a
question with an answer. A cracked slab is a hazard to one person and an
inconvenience to another; whether a corner \emph{needs} a curb ramp depends on
what counts as a pedestrian route. These are judgement calls, and cities now
make them at scale by asking volunteers.

Project Sidewalk~\cite{saha2019projectsidewalk} is among the largest such
efforts. Volunteers examine street-level imagery and mark accessibility
problems---missing curb ramps, surface defects, obstructions---and other
volunteers review each mark as \emph{agree}, \emph{disagree}, or
\emph{not sure}. The resulting labels feed accessibility audits, routing for
wheelchair users, and municipal remediation priorities. Before any of that, the
data is cleaned: labels that failed to attract sufficient agreement are filtered
out~\cite{li2024labelaid}.

That filtering step embeds an assumption worth examining. It treats a label
reviewers argued about as defective, as though disagreement indicates a mistake
by the original labeller. But disagreement among careful reviewers can also
indicate that the underlying question is genuinely contested, which is a
different thing and carries different consequences. If reviewers split over
whether a particular crack rises to the level of a surface problem, that split
is evidence about where a city's accessibility threshold actually sits and how
widely it is shared. Deleting it does not resolve the ambiguity; it conceals it.

Disagreement is also not one phenomenon. A reviewer can decline to endorse a
label in two ways that look similar in aggregate and mean opposite things.
Reviewers can \emph{split}: some confidently agree, others confidently
disagree, and the argument is about where the standard lies. Or reviewers can be
\emph{unsure}: nobody takes a position, because the image does not settle the
question. The first is a disagreement about values, which more reviewing will
not converge. The second is a disagreement about evidence, which better imagery
could in principle fix. A pipeline that cannot tell them apart cannot route
either correctly.

Project Sidewalk is unusual in making the distinction separable. Most
crowdsourcing interfaces offer only agreement or disagreement; the explicit
``not sure'' option, recorded per reviewer alongside reviewer identity and
position in the review sequence, is what makes the decomposition possible.

We make three contributions:

\begin{itemize}
\setlength\itemsep{0.15em}
\item \textbf{An audit of what quality control removes.} The published filter
  retains over half of all reviewed labels but under a tenth of contested ones,
  in every one of twelve cities, and the removal is spatially uneven by a
  factor of three to eight beyond sampling noise
  (Sections~\ref{sec:qc} and~\ref{sec:spatial}).
\item \textbf{Evidence that filtering manufactures overconfidence.} A model
  trained on the filtered record is far more certain about exactly the labels
  reviewers could not agree on, at no gain in accuracy, across three model
  families (Section~\ref{sec:calibration}).
\item \textbf{An empirical decomposition of the disagreement itself} into a
  standards component and an evidence component, with different sources,
  corroborated independently by the labellers' own tags, and neither of which
  resolves with additional review (Sections~\ref{sec:twokinds}--\ref{sec:resolve}).
\end{itemize}

A methodological thread runs alongside these. Each finding changed sign or
magnitude when we controlled properly for who reviewed a label and which city it
came from. We report the uncontrolled analyses beside the controlled ones,
because in every case the uncontrolled version told a tidier story than the data
supports, and because anyone repeating this work on similar data will meet the
same trap.

\section{Related Work}

\textbf{Disagreement as signal.} A substantial literature argues that annotator
disagreement carries information rather than noise. \citet{aroyo2015truth} frame
inter-annotator disagreement as intrinsic to semantic annotation, and CrowdTruth
\cite{dumitrache2018crowdtruth} operationalises this with metrics that model the
interdependence of worker, input unit, and annotation---separating worker
quality from unit ambiguity. Our label-versus-reviewer decomposition is an
instance of that idea; our contribution is not the decomposition but what it
reveals when applied to a deployed civic pipeline.

\textbf{Learning from disagreement.} Training on aggregated hard labels is known
to produce overconfident models, and preserving the label distribution improves
calibration \cite{peterson2019human, uma2021learning}. Our calibration result
applies this established effect to a real cleaning pipeline and measures its
magnitude in a setting with policy consequences, rather than proposing a new
learning method.

\textbf{Bias in volunteered geographic information.} Coverage bias in
crowdsourced spatial data is well documented: contributions concentrate in
particular places and are shaped by who contributes \cite{haklay2010good,
quattrone2015bias}. That literature concerns \emph{where volunteers choose to
map}. We examine a different and less-studied mechanism: bias introduced at the
\emph{retention} step, after data has been collected, by the filter that decides
what survives.

\textbf{Crowdsourced accessibility assessment.} Project Sidewalk
\cite{saha2019projectsidewalk} and subsequent work on label quality
\cite{li2024labelaid} establish the platform and its validation pipeline. We
take that pipeline as our object of study.

\section{Data}
\label{sec:data}

We use every label from the twelve live Project Sidewalk city deployments,
retrieved through the platform's public API in August 2026: 811{,}368 labels in
total, of which 177{,}599 carry at least three independent reviews. Those
reviewed labels rest on 746{,}975 individual review decisions contributed by
11{,}480 reviewers. Each label carries its type, a severity rating, optional
tags and free text, geographic coordinates and neighbourhood, the labeller's
identifier, and the full sequence of reviews with each reviewer's identity.

We decompose each label's review record into two quantities. \emph{Split}
measures how evenly the reviewers who took a position divided between agree and
disagree, and is $1.0$ when they divided equally and $0$ when unanimous.
\emph{Unsure} is the share of reviewers who declined to take a position. A label
can score high on either, both, or neither. Where a categorical grouping is
needed we call a label \emph{split} when its split score exceeds $0.6$ with
fewer than 20\% unsure votes, and \emph{unsure} when at least a third of its
reviewers answered ``not sure''.

\section{Two Kinds of Disagreement}
\label{sec:twokinds}

Both populations are substantial. Reviewers split into opposing camps on
23{,}638 labels (13.3\% of the reviewed record) and answered ``not sure'' at a
rate of one third or more on 7{,}299 (4.1\%). The two measures are only weakly
related ($\rho = 0.14$), the first indication that they are distinct phenomena
rather than two readings of a single underlying difficulty.

Disagreement is not uniform across the label taxonomy (Table~\ref{tab:types}).
Missing curb ramps and signal-related labels draw the highest rates of splitting,
consistent with judgements about whether an intervention \emph{ought} to exist
rather than about what is visible. Obstacles and surface problems draw the
highest rates of ``not sure''. Splitting also rises with the severity rating the
original labeller assigned, from 10.9\% at severity 1 to 15.5\% at severity 2,
which is what a disagreement about where a threshold sits would predict.

\begin{table}[t]
\caption{Disagreement by label type, twelve cities. Types differ in both the
rate and the kind of disagreement they attract.}
\label{tab:types}
\begin{tabular}{lrrrr}
\toprule
\textbf{Label type} & \textbf{$n$} & \textbf{Split} & \textbf{Unsure} & \textbf{QC keeps} \\
\midrule
Crosswalk       & 14{,}717 &  6.2\% & 0.9\% & 74.8\% \\
CurbRamp        & 54{,}676 & 10.3\% & 1.8\% & 71.9\% \\
NoCurbRamp      & 32{,}371 & 21.4\% & 2.8\% & 52.7\% \\
Obstacle        & 26{,}102 & 12.5\% & 8.9\% & 31.1\% \\
Signal          &  4{,}293 & 24.6\% & 1.1\% & 68.7\% \\
SurfaceProblem  & 45{,}221 & 12.6\% & 6.4\% & 40.4\% \\
\bottomrule
\end{tabular}
\end{table}

\subsection{Corroboration from the labellers' tags}

The decomposition receives independent support from a field it never touched.
Labellers may attach descriptive tags to a mark. Among split labels, the most
frequent informative tags concern routing: \emph{no alternate route} (2{,}490
labels) and \emph{alternate route present} (1{,}595). Among unsure labels, the
most frequent concern visibility and occlusion: \emph{pole} (409),
\emph{narrow sidewalk} (398), \emph{parked car} (306).

Reviewers argue about whether a barrier \emph{matters}, and they are unsure when
something \emph{blocks the view}. That the two populations attract different
vocabulary, written by the original labellers rather than by the reviewers whose
votes define the split, is stronger evidence for the distinction than any
statistic computed from the votes themselves.

\section{Where the Disagreement Lives}
\label{sec:where}

Reviewers are not interchangeable. Among the 3{,}333 reviewers with at least
twenty reviews, ``not sure'' rates range from 0\% to 70\% with a median of
3.0\%. A raw count of unsure votes on a label therefore confounds a property of
the sidewalk with a property of whoever happened to review it.

We separate the two by fitting an additive model
$y_{ij} \sim \mu + a_i + b_j$ over individual review decisions, where $a_i$ is a
label effect and $b_j$ a reviewer effect, estimated by alternating means and
shrunk toward zero in proportion to the number of reviews supporting each label.

The model must be fitted \emph{within} each city. Deployments differ in both
their reviewer populations and their base rates, so a pooled fit absorbs
between-city variation into the reviewer term and overstates it: the pooled
estimate attributes 20.8\% of the variance in ``not sure'' responses to
reviewers against 14.3\% to labels, whereas within-city estimates reverse that
ordering in ten of twelve deployments.

Fitted correctly, both outcomes are label-driven and differ in degree rather
than kind (Table~\ref{tab:variance}). Splitting is a property of the label,
exceeding the reviewer term by roughly five to one and doing so in all twelve
cities. Uncertainty is only weakly so: reviewer identity accounts for about
twice as much of the variance in ``not sure'' as in disagreement.

\begin{table}[t]
\caption{Variance in individual review decisions attributable to the label
versus the reviewer, estimated within each city (median over twelve
deployments).}
\label{tab:variance}
\begin{tabular}{lcc}
\toprule
\textbf{Source} & \textbf{Disagree} & \textbf{Unsure} \\
\midrule
Label            & \textbf{52.0\%} & 27.1\% \\
Reviewer         & 9.7\%           & 20.7\% \\
Label:reviewer   & \textbf{5.3:1}  & 1.3:1 \\
Label-dominated in & 12/12 cities  & 10/12 cities \\
\bottomrule
\end{tabular}
\end{table}

The practical consequence is that raw and reviewer-adjusted estimates of which
labels are unclear correlate at only $\rho = 0.68$: removing reviewer tendency
substantially reorders the population.

\section{More Review Does Not Produce Consensus}
\label{sec:resolve}

The review records preserve each reviewer's identity and position within the
label's review sequence. That position carries real ordering rather than being
arbitrary: shuffling positions within each label collapses the between-reviewer
variance in mean position from $0.024$ to $0.004$ ($z=195$), and reviewer mean
positions span $0.00$ to $0.87$ where an arbitrary order would place every
reviewer near $0.5$. The ordering is not alphabetical by reviewer identifier
(6.2\% of labels are so ordered, against 6.7\% expected by chance), and
reviewers who joined later occupy later positions ($\rho = 0.29$,
$p < 10^{-26}$, using label creation order as an arrival proxy). We therefore
treat position as sequence.

Given that, we compare a label's first three reviews against its later ones
\emph{within the same label}, which removes any selection effect arising from
which labels attract additional review. We restrict to labels with at least six
reviews and subtract each reviewer's estimated tendency, fitted within their own
city.

The result is a null, and it is the point. Neither form of disagreement declines
over a label's lifetime. Uncertainty rises slightly and consistently: the
adjusted delta is positive in eight of eleven deployments with sufficient data,
median $+0.008$. Disagreement shows no consistent direction, positive in three
of eleven, median $-0.005$.

Two analyses we report but do not rely on illustrate why the controls matter.
Without adjusting for reviewer tendency, uncertainty appears to \emph{decline}
over a label's lifetime ($0.106 \rightarrow 0.096$); that apparent resolution is
explained by early reviewers being more unsure-prone than later ones (mean
reviewer effect $+0.043$ versus $+0.020$). And when reviewer effects are pooled
across deployments rather than fitted within them, disagreement appears to rise
sharply ($+0.026$); that disappears once each city's reviewers are modelled
separately. Both artefacts point toward a tidier story than the data supports.

\section{Quality Control Removes the Contested Cases}
\label{sec:qc}

Project Sidewalk's published quality criterion retains a label when it has more
than two agreements and no more than two disagreements. Applied to our corpus it
retains 54.5\% of reviewed labels overall (95\% CI [54.2, 54.7]). Applied to the
contested populations it retains 8.6\% of split labels ([8.3, 9.0]) and 6.8\% of
unsure ones ([6.2, 7.3])---discarding 91\% and 93\% respectively.

This is not an artefact of any one deployment. Contested labels are retained
below the city average in all twelve cities, and leaving any single city out
moves the pooled figures by at most 1.3 percentage points; the largest shift, to
9.9\% and 7.5\%, comes from excluding Seattle, which alone contributes 30\% of
the reviewed corpus.

The criterion also has a property worth stating plainly: because it requires
more agreements than disagreements, every label it retains is one the majority
endorsed. All of them are, by construction, ``correct''. The filter cannot
define a training set on its own, and a realistic pipeline must pair it with a
mirror rule selecting consensus negatives. Under that pairing, 54.5\% of the
reviewed record is consensus-positive, 8.9\% consensus-negative, and
\textbf{36.6\% is an ambiguous middle that is discarded entirely}.

\section{Removal Is Spatially Uneven}
\label{sec:spatial}

Contested labels are not distributed uniformly across a city, so the filter
removes evidence unevenly. Across 284 neighbourhood regions with at least 100
reviewed labels, between-region variation in retention exceeds binomial
sampling noise by a factor of 3.3 to 7.6 in every city. Seattle spans 11\% to
71\% retention across its 83 regions; Chicago spans 20\% to 83\% across 78.

The obvious explanation---that regions differ in what gets labelled there, and
some label types are more contested---does not account for it. Holding label
type composition fixed leaves spreads of 15 to 57 percentage points, and 68 of
284 regions retain more than ten points below their type-expected rate.

The consequence is that the accessibility map produced after cleaning is
systematically more complete in some neighbourhoods than others. A city
allocating remediation budget from that map sees fewer problems where labels
were contested---not because there are fewer, but because the record was thinner
after cleaning.

\section{Filtering Manufactures Overconfidence}
\label{sec:calibration}

Downstream models are trained on what survives. If removing contested cases also
removes a model's opportunity to learn that some cases are genuinely uncertain,
then routine cleaning produces a system that is confident precisely where people
disagree.

We test this without touching an image. We predict whether a label will be
judged correct from 39 features available at labelling time---label type,
severity, tags, geometry, imagery age, and the labeller's leave-one-out
history---excluding all review counts, which constitute the target. Folds are
split by neighbourhood so no region appears in both training and test, and every
regime is evaluated on the same held-out set. We compare four training regimes:
\emph{filtered}, using only consensus labels as a real pipeline would;
\emph{full\_sub}, subsampled to the filtered set's size but retaining the
ambiguous middle; \emph{full}, all reviewed labels; and \emph{soft}, all labels trained against
the empirical review distribution. For \emph{soft} we enter each label twice,
as a positive weighted by its agreement rate and a negative weighted by the
remainder, which is cross-entropy against the review distribution and keeps the
model class identical to the other three regimes.

\begin{table}[t]
\caption{Predicting label validity under four training regimes, evaluated on the
same held-out labels. ``Silent'' is the fraction of contested test labels the
model calls with $p > 0.9$.}
\label{tab:calibration}
\begin{tabular}{lrrrr}
\toprule
\textbf{Regime} & \textbf{AUROC} & \textbf{ECE} & \textbf{Conf.} & \textbf{Silent} \\
\midrule
filtered  & 0.900 & 0.066 & 0.703 & \textbf{48.0\%} \\
full\_sub & 0.921 & 0.030 & 0.600 & 30.3\% \\
full      & 0.923 & 0.030 & 0.603 & 29.8\% \\
soft      & 0.907 & 0.084 & 0.484 & \textbf{14.4\%} \\
\bottomrule
\end{tabular}
\end{table}

Table~\ref{tab:calibration} gives the result. Training on the filtered record
makes the model assert more than 90\% confidence on 48.0\% of the labels
reviewers split on, against 30.3\% when the ambiguous middle is retained at
matched sample size. Sample size therefore explains none of the gap; composition
explains all of it.

There is no accuracy argument for the filtering. AUROC is \emph{lower} for the
filtered model (0.900 against 0.921), and its calibration error is twice as
large. The filter costs calibration and buys nothing measurable.

The effect is not specific to one learner. Repeating the matched comparison with
a random forest and with logistic regression gives gaps of $+21.9$ and $+30.4$
percentage points respectively, against $+17.9$ for gradient boosting, and the
filtered model has lower AUROC in all three.

\subsection{A remedy that costs nothing}

The same table contains the fix. Training against the review distribution rather
than a hard majority reduces silent assertions from 48.0\% to 14.4\%, a
3.3-fold reduction, while AUROC \emph{rises} slightly (0.900 to 0.907). The
model is not trading accuracy for humility; it is declining to invent certainty
the reviewers never expressed.

The change is small and requires no new data. The review counts already exist in
the database---they are simply discarded during aggregation. Retaining a label
with the note that eleven reviewers agreed and nine did not, rather than
recording it as ``correct'', is enough.

We note one metric artefact. Expected calibration error is \emph{worse} for the
soft regime (0.084 against 0.066), because ECE is computed against hard majority
labels and therefore rewards confidence on targets that are close to coin flips
by construction. On contested labels the majority verdict carries little
information, and a model penalised for hedging on them is being penalised for
the behaviour we want. Where the evaluation target is itself contested, ECE
against a thresholded majority is the wrong instrument.

A regressor trained on the agreement rate directly reaches 12.9\%, confirming
that the reduction is not an artefact of how soft targets are encoded.

\section{Discussion}

\textbf{Route, do not delete.} The two kinds of disagreement have different
remedies. Labels where reviewers were unsure describe a deficit of evidence:
better imagery, a different capture angle, or a site visit could settle them.
Labels where reviewers split describe a deficit of shared standards, and no
amount of additional review will converge them---as Section~\ref{sec:resolve}
shows empirically. Those belong in front of a person with the authority to
decide what the city's threshold is. A pipeline that discards both treats a
policy question as a data-quality problem.

\textbf{Contested labels are a municipal asset.} A city that knows which of its
accessibility judgements are contested knows where its standards are unclear.
That is actionable information, and it is currently being deleted before anyone
can read it. Preserving it costs nothing: the review records already exist.

\textbf{Cleaning has spatial consequences.} Because contestation clusters,
uniform filtering produces non-uniform coverage. Any audit of a crowdsourced
accessibility map should report retention by neighbourhood alongside coverage,
since a map can be complete in collection and incomplete after cleaning.

\textbf{Aggregation, not filtering, is the deeper issue.} Our filtered-versus-%
matched comparison isolates the cost of \emph{deleting} contested labels, and
our soft-target comparison suggests a second cost of the same order from
\emph{collapsing} the retained ones to hard labels. Even a pipeline that deleted
nothing would discard the agreement margin, presenting a 20--1 label and a 6--4
label as equally certain.

\section{Limitations}

Our downstream model predicts label validity from metadata rather than from
imagery, so it is a proxy for the computer-vision models the platform actually
trains \cite{li2024labelaid}. Whether the calibration effect transfers at the
same magnitude to an image model is untested and is the clearest next step.

The overconfidence result is an application of a known effect---hard labels
produce overconfident models---to a deployed pipeline. Its contribution is the
audit and the magnitudes, not a new claim about learning.

Contested labels have near-arbitrary majority verdicts by construction: for a
17--16 split, the target is close to a coin flip. Confidence on such labels is
therefore confidence about a noisy target. We regard that as the point rather
than a confound, since the model has no access to the review counts, but a
reader may weigh it differently.

Our sequence analysis rests on review position reflecting arrival order.
Section~\ref{sec:resolve} gives three empirical tests supporting this, but we
have not confirmed with the platform maintainers that the ordering is
specifically chronological rather than database insertion order, though the two
normally coincide.

Finally, the additive variance model is a simplification; a crossed
random-effects logistic fit would be the stricter treatment.

\section{Conclusion}

Quality control in crowdsourced urban accessibility data removes disagreement
rather than error. Over 90\% of the labels reviewers contested are discarded in
every one of twelve cities, the removal falls unevenly across neighbourhoods,
and models trained on what survives become confident about precisely the cases
people could not agree on---at no gain in accuracy. The disagreement being
deleted is not noise about the sidewalk; it is a record of where a city's
accessibility standards are unsettled. It should be routed to the people who can
settle them, not filtered away before anyone sees it.

\bibliographystyle{ACM-Reference-Format}
\bibliography{refs}

\end{document}
